\subsection{Additional Details}
\label{sec:deets}

We used 10-fold cross validation to evaluate the performance (fairness
and accuracy losses) of a trained model on test data. For smaller
datasets (Communities and Crime, COMPAS and Sentencing) we used one
run of 10-fold cross validation. For larger datasets (Adult, Default
and Law School) we first randomly sampled 30-50\% of the dataset
(depending on the dataset), ran the 10-fold cross validation on the
sampled data and then repeated the experiment 3 times each time with a
new random sample. This is because both the running time and the
instability of the CVXPY solver would increase when we used the whole
datasets for the larger datasets.

While the fairness losses~\ref{eqn:innersq}, \ref{eqn:outersq}, and
\ref{eqn:betweensq} are defined using all the $n_1\times n_2$ cross
pairs in the dataset, in our experiments we only used
$2 \times \text{Minority } n$ random cross pairs where
$\text{Minority } n = \min\{n_1, n_2\}$ (see
Table~\ref{tab:summary}). This is because: (1) using more cross pairs
did not substantially improve the efficiency curves in
Figure~\ref{fig:pareto}, (2) the CVXPY solver for binary-valued
problems would become unstable when using individual fairness if we
increase the number of cross pairs significantly. Furthermore, we used
the same cross validation folds and cross pairs when experimenting
with our different notions of fairness.

Finally, linear models perform poorly on some of our datasets~(see
e.g.~\cite{default-dataset}). However, our goal in this work was not
to obtain the most accurate models for particular datasets but to
study the trade-off between fairness and accuracy and also different
notions of fairness for regression problems.
